\section{INTRODUÇÃO}

Nos últimos anos, a Engenharia de Software passou por uma transformação significativa com a consolidação da cultura DevOps e a adoção massiva da computação em nuvem. Nesse cenário, a prática de Infraestrutura como Código (IaC) emergiu como o padrão para provisionar e gerenciar recursos de forma escalável, auditável e versionável, utilizando ferramentas declarativas de mercado, como o Terraform. O presente projeto de pesquisa tem como objeto de estudo a interseção entre essas práticas avançadas de automação de infraestrutura e a utilização emergente de modelos de Inteligência Artificial Generativa.

Apesar dos benefícios inegáveis da IaC, a complexidade crescente dos ecossistemas de nuvem introduziu novos desafios operacionais. A gênese do problema investigado reside quando ocorrem falhas nas esteiras de entrega contínua (CI/CD) durante o provisionamento de recursos. Os logs de erro gerados por falhas de sintaxe, dependências cíclicas ou bloqueios de permissão na nuvem são frequentemente extensos, prolixos e não lineares. Consequentemente, profissionais de operações e engenheiros de confiabilidade (SRE) despendem um tempo excessivo no troubleshooting manual, cruzando dados de logs fragmentados com extensas documentações técnicas. Esse gargalo humano eleva o Tempo Médio de Recuperação (MTTR) dos sistemas e expõe gargalos na eficiência das equipes.

A literatura recente sobre AIOps (Inteligência Artificial para Operações de TI) tem demonstrado que Modelos de Linguagem de Grande Escala (LLMs) possuem notável capacidade para processar dados semiestruturados e auxiliar na geração e correção de código. No entanto, o uso empírico, isolado e desestruturado dessas ferramentas na análise de falhas críticas esbarra em limitações conhecidas, como as ``alucinações'' --- situações em que a IA sugere configurações incorretas ou inseguras por falta de contexto arquitetural adequado. Autores e pesquisadores da área de engenharia de software apontam a necessidade de se estabelecer parâmetros mais rígidos e metodologias que guiem a interação entre o engenheiro e a IA garantindo que as sugestões geradas sejam rápidas, precisas e, acima de tudo, aderentes às boas práticas de segurança corporativa.

É neste contexto que a presente pesquisa se insere, buscando preencher a lacuna entre o uso ad-hoc de inteligência artificial em fóruns e chats e a necessidade de processos padronizados em ambientes corporativos de alta criticidade.

Diante do exposto, o presente trabalho é norteado pela seguinte pergunta de pesquisa: \textbf{Como um método estruturado de análise de logs via IA Generativa pode auxiliar na identificação e proposição de correções para falhas de provisionamento em ferramentas de infraestrutura como código?} Com o intuito de responder a este questionamento, o projeto tem como objetivo principal propor um fluxo de trabalho metodológico e arquitetural que otimize a injeção de contexto na IA e parametrize a leitura de logs, garantindo maior resiliência, assertividade e agilidade na tomada de decisão das equipes de DevOps.
